\documentclass[11pt,a4paper]{article}
\usepackage[T1]{fontenc}
\usepackage[utf8]{inputenc}
\usepackage{lmodern}
\usepackage[margin=2cm]{geometry}
\usepackage{booktabs}
\usepackage{array}
\usepackage{longtable}
\usepackage{xcolor}
\usepackage{enumitem}
\usepackage{titlesec}
\usepackage{hyperref}
\usepackage{fancyhdr}
% Allow long identifiers in \texttt to hyphenate instead of running into
% the right margin.
\usepackage[htt]{hyphenat}
\setlength{\emergencystretch}{3em}

\definecolor{osdag}{HTML}{116A5B}
\hypersetup{colorlinks=true, linkcolor=osdag, urlcolor=osdag}
\titleformat{\section}{\large\bfseries\color{osdag}}{\thesection}{0.6em}{}
\titleformat{\subsection}{\normalsize\bfseries}{\thesubsection}{0.5em}{}
\setlist{itemsep=2pt,topsep=3pt,parsep=0pt}
\pagestyle{fancy}
\fancyhf{}
\fancyhead[L]{\small OsdagBridge --- LaTeX Report Generator Refactoring}
\fancyhead[R]{\small FOSSEE Winter Internship 2026}
\fancyfoot[C]{\small\thepage}
\renewcommand{\headrulewidth}{0.4pt}

\newcommand{\code}[1]{\texttt{\small #1}}

\title{\vspace{-1.2cm}\textbf{Technical Changes Document}\\[2pt]
\large OsdagBridge LaTeX Report Generator --- Refactoring \& Enhancements}
\author{Screening Task Submission --- FOSSEE Winter Internship 2026}
\date{}

\begin{document}
\maketitle
\vspace{-0.6cm}

\section{Objective and approach}

The OsdagBridge report generator produces a multi-chapter structural design report
(design basis, loads, analysis, member checks, BOQ) by assembling LaTeX source from
Python string builders. The task asked for six improvements: repeated table headers,
layout/footer fixes, cleaner load and geometry tables, a utilization-ratio chart,
material-quantity charts, and a centralized formatting module.

All six were implemented on the production code path
(\code{build\_report\_payload()} $\rightarrow$ \code{generate\_report()}), not on a
side branch of logic. Because the design engine depends on \code{openseespy}, which
has no wheel for the Python 3.14 interpreter available on this machine, the demo
PDFs are produced by an inert OpenSees stub
(\path{tools/report_demo/generate_demo_report.py}) and a fixed 30\,m four-girder
dataset (\path{tools/report_demo/demo_data.py}). The same payload drives both the
before and after PDFs, so every difference between them is attributable to the code
changes and not to different input data.

The \emph{before} PDF is generated from an untouched checkout of commit
\code{d74f56d} (\code{git worktree}), using the identical dataset.

\section{Files changed}

\begin{longtable}{@{}>{\raggedright\arraybackslash}p{6.0cm}>{\raggedright\arraybackslash}p{9.8cm}@{}}
\toprule
\textbf{File} & \textbf{Change} \\
\midrule
\endhead
\code{reports/theme.py} & \textbf{New.} Single source of truth: palette, page
geometry, header/footer, table style, \code{preamble()}, \code{longtable()},
\code{figure()}, header-repetition and page-break normalisation passes,
\code{apply\_chart\_theme()}. \\
\code{reports/charts.py} & \textbf{New.} Matplotlib generators for the
utilization-ratio chart and the material-quantity charts. \\
\code{reports/report\_generator.py} & Preamble delegated to \code{theme.preamble()};
legacy preamble removed; document-level post-processing applied; chart asset
directory wired into Chapters 5 and 7. \\
\code{reports/chap2.py}, \code{chap4.py}, \code{chap8.py} & Hard-coded table
padding and row-stretch commands removed; wide tables now call
\code{theme.table\_padding()}. \\
\code{reports/chap3.py} & Live-load table split into a vehicle table and a separate
footway table, with explicit units and UI-selected vehicles. \\
\code{reports/chap5.py} & Utilization ratios collected while the check tables are
built and rendered as a summary chart in Section~5.5. \\
\code{reports/chap7.py} & Rewritten on \code{theme.longtable()}; material-quantity
charts appended to the chapter. \\
\code{tools/report\_demo/} & \textbf{New.} Headless runner and demo dataset used to
produce the submitted PDFs. \\
\bottomrule
\end{longtable}

\section{1. Repeated table headers across page breaks}

\textbf{Before.} Of the 49 longtables in the report, only 2 declared
\code{\textbackslash endhead}; a table that crossed a page boundary continued
without any column headings, so values on the new page could not be mapped to
columns.

\textbf{After.} All 51 longtables now carry a head block, and continuation pages
repeat the header plus the caption with a ``(continued)'' suffix.

Two mechanisms were used deliberately:
\begin{itemize}
  \item \code{theme.longtable()} emits \code{\textbackslash endfirsthead},
        \code{\textbackslash endhead} and \code{\textbackslash endlastfoot} blocks
        for every table it builds (Chapters 3, 5, 7).
  \item \code{theme.repeat\_longtable\_headers()} is a document-level pass applied in
        \code{generate\_report()}. It parses each \code{\textbackslash begin\{longtable\}}
        block, extracts the column specifier and trailing rule, and injects the head
        blocks after the first \code{\textbackslash endfirsthead} or
        \code{\textbackslash endhead}. This fixes the pre-existing hand-written tables
        in Chapters 2, 4, 6 and 8 without touching each of them individually.
\end{itemize}

\section{2. Page geometry, vertical spacing and footer overlap}

\textbf{Before.} The report used \code{geometry} with 1\,in margins, a 15\,pt
\code{headheight} (LaTeX warned that it was too small) and no page-bottom guard
before tables. Table 2.x and 5.1 could start with two rows left at the bottom of a
page, and a group label centred with \code{\textbackslash multirow} could be pushed
into the footer margin --- the reported ``G4 bleeding into the footer'' defect.

\textbf{After.}
\begin{itemize}
  \item Geometry is defined once in \code{theme.PageGeometry} (A4, 2\,cm margins,
        18\,pt head height, \code{includeheadfoot}) and emitted by
        \code{theme.preamble()}.
  \item \code{theme.TableStyle} carries \code{needspace\_baselines = 7}; the global guards
        were removed (they produced a stray ``(continued)'' head at the page bottom).
  \item Cell padding is a theme constant (\code{tabcolsep = 3.0pt}); unusually wide
        tables request \code{theme.table\_padding(colsep=TIGHT\_COLSEP)} instead of
        hard-coding their own values.
  \item \code{theme.flatten\_breakable\_multirows()} replaces each
        \code{\textbackslash multirow} group-label cell with a plain
        \code{\textbackslash makecell} on the group's first row. The
        existing \code{\textbackslash cline} rules preserve the visual grouping, but
        no label box can any longer be vertically centred across a page break, so no
        label can be drawn below the bottom margin. This is applied centrally, so it
        also protects tables added later.
\end{itemize}

\section{3. Load and geometry tables}

\textbf{Before.} Vehicular live load, impact factor and footway loading were merged
into a single Live Loads table with mixed units in one column; the footway
contribution was not separately verifiable.

\textbf{After.} Chapter 3 now contains:
\begin{itemize}
  \item Table 3.3 --- \emph{Vehicular Live Load (IRC 6)}, one row per vehicle selected
        in the UI (Class A, Class 70R, etc.), with impact factor and IRC reference;
  \item Table 3.4 --- \emph{Governing Longitudinal Load Effects}, with braking and
        centrifugal forces and their IRC references;
  \item Table 3.5 --- \emph{Footway Live Load}, with its own value/unit columns and
        the IRC clause reference for each parameter;
  \item Table 3.6 --- wind load, unchanged in content but rebuilt on the theme table
        builder.
\end{itemize}
Units are now a dedicated column in both tables, so a number is never ambiguous.
Chapter 3 reads the selected vehicles from the same payload keys the GUI writes, so
the table always matches what the user ticked in the Loading tab.

\section{4. Utilization-ratio chart}

While the detailed check tables are generated, \code{chap5.py} records the numeric
demand/capacity ratio for each check into \code{\_ur\_values}. Section 5.5
(\emph{Overall Design Check Summary}) then renders a four-category bar chart ---
Steel Plate Girders, Concrete Deck Slab, Cross Bracing, End Diaphragms --- with:
\begin{itemize}
  \item UR on the vertical axis and a dashed red line at UR = 1.0 (the limit state);
  \item value labels on each bar and colour coding (green for UR $\leq$ 1.0,
        red above 1.0);
  \item the same palette and font as the rest of the report, applied through
        \code{theme.apply\_chart\_theme()}.
\end{itemize}
The figure is written to the temporary asset directory and embedded through
\code{theme.figure()}, so the report remains self-contained when the temporary
directory is cleaned up.

\section{5. Material quantity charts}

Chapter 7 was rewritten on top of \code{theme.longtable()} and now ends with a
\emph{Material Quantity Visualisation} section containing two charts:
\begin{itemize}
  \item steel tonnage per element (girders, cross bracing, end diaphragms, shear
        studs) --- only elements with a calculated tonnage are plotted;
  \item deck concrete volume ($\mathrm{m}^3$) and reinforcement steel weight (MT)
        as a paired bar chart.
\end{itemize}
Both are generated at report time into the temporary asset directory and embedded
as PNG, so no manual chart files are required. Because the BOQ calculation is the
single source of the plotted numbers, the charts and Table 7.1 always agree.

\section{6. Centralized formatting module}

\code{theme.py} is now the only place that knows about visual style. It exposes:
\begin{itemize}
  \item \code{ReportStyle} / \code{PageGeometry} / \code{TableStyle} /
        \code{HeaderFooterStyle} dataclasses with the palette, fonts, rules and
        spacing;
  \item builders \code{preamble()}, \code{longtable()}, \code{figure()},
        \code{table\_padding()};
  \item document passes \code{repeat\_longtable\_headers()} and
        \code{flatten\_breakable\_multirows()};
  \item \code{apply\_chart\_theme()}, which configures Matplotlib so figures use the
        report's fonts and colours.
\end{itemize}
Chapter modules now import formatting from \code{theme} and contain only content.
All per-chapter \code{\textbackslash setlength}, \code{\textbackslash renewcommand}
and \code{\textbackslash captionsetup} calls were removed or redirected through the
theme helpers.

\section{Verification}

Both PDFs were compiled with the same TeX Live installation (TeX Live 2026,
\code{pdflatex}, three passes) from the same dataset.

\begin{center}
\begin{tabular}{@{}lrr@{}}
\toprule
\textbf{Metric} & \textbf{Before} & \textbf{After} \\
\midrule
LaTeX errors (\code{!} lines in log)          & 138 & 0 \\
Overfull box warnings (hbox + vbox)            & 61  & 0 \\
Pages (incl. unnumbered title page)            & 48  & 47 \\
Longtables with a repeating head               & 2 / 49 & 51 / 51 \\
Tables continuing headerless                   & yes & no \\
Content overlapping the footer box             & yes & 0 pages (bbox scan) \\
Utilization-ratio chart                        & absent & present (Sec.~5.5, 4 elements) \\
Material-quantity charts                       & absent & present (Ch.~7) \\
Vehicle LL unit column                         & absent & present \\
Per-chapter hard-coded spacing commands        & 6 files & 0 files \\
\bottomrule
\end{tabular}
\end{center}

The after build is clean: zero overfull boxes (horizontal or vertical) and zero
LaTeX errors. The 138 errors in the before run were fatal
table-construction errors (\code{Misplaced \textbackslash omit},
\code{Misplaced \textbackslash noalign}, \code{Extra alignment tab}) caused by
\code{\textbackslash caption} rows that were not terminated with
\code{\textbackslash\textbackslash}; those rows are now emitted by
\code{theme.longtable()} with correct termination.

\section{How the PDFs were produced}

\begin{verbatim}
# after (this branch)
python tools/report_demo/generate_demo_report.py --out out --stem Report_After

# before (pristine checkout of d74f56d, same dataset)
git worktree add ../baseline d74f56d
cp -r tools/report_demo ../baseline/tools/
python ../baseline/tools/report_demo/generate_demo_report.py \
       --out out --stem Report_Before
\end{verbatim}

Each run writes both the \code{.tex} source and the compiled \code{.pdf}. The runner
reports LaTeX errors by scanning the log, so a regression fails loudly instead of
producing a silently broken PDF.

\section{Additional defects found and fixed}

While verifying the six requirements, four further defects surfaced. They are
listed separately because the task did not name them.

\subsection*{Braking force always reported as ``N/A''}

\code{IRC6\_2017.cl\_211\_2\_braking\_force()} referenced
\code{braking\_force\_2} before assigning it whenever the bridge had one or two
traffic lanes, raising \code{UnboundLocalError}; the report caught the exception
and printed ``N/A''. The function also returned newtons labelled as tonnes. Both
are fixed: the terms are initialised to zero, wheel loads (stored in newtons) are
converted to tonnes, and additional lanes accumulate their 5\,\% share. The
report now shows 108.69\,kN for the two-lane demo bridge. The function had no
other callers, so no design result changes.

\subsection*{``Page X of Y'' footer lagging by one}

The \code{lastpage} label is written when the final page is shipped, so it
trapped the page count as it was before the closing page break and printed
``of 50'' in a 51-page file. The generator now runs three passes and the label
is written explicitly on the true last page. The references page also uses
\code{\textbackslash chapter*}, which silently switches to the plain page style;
it now restores the report page style, so the final page carries the same header
and footer as the rest of the document.

\subsection*{Chapter 7 quantities were placeholders}

The BOQ returned ``N.A.'' for every steel item because the demo dataset lacked
the section keys the calculator reads. The dataset was completed, and end
diaphragm tonnage --- previously not computed at all --- is now derived in
\code{boq\_generator.py} from the transverse-bracing section and geometry
(two frames per bay, two diagonals and two chords each). The Chapter 7 table and
the steel chart therefore show the same, computed number.

\subsection*{Utilization summary inflated by slenderness}

The KL/r check was being averaged into the cross-bracing utilization ratio even
though it is not a demand/capacity ratio. Slenderness rows are now excluded from
the chart while remaining in the detailed table.

\section{Structure of the central style module}

\begin{verbatim}
theme.py
|-- PALETTE, SERIES_COLORS, apply_chart_theme()
|-- PageGeometry      paper, margins, head height, footskip
|-- TableStyle        font size, colsep, arraystretch, rule widths,
|                     needspace threshold, caption placement
|-- HeaderFooterStyle left/right header text, footer text, rule colour
|-- ReportStyle       font size, line stretch, title/heading fonts
|-- preamble()        emits the complete LaTeX preamble from the above
|-- longtable()       colspec + head/body/foot blocks + caption + note
|-- figure()          [H] float, caption, width, centring
|-- table_padding()   tight-padding helper for wide numeric tables
|-- repeat_longtable_headers()   document-level header pass
|-- flatten_breakable_multirows() page-break-safe group labels
`-- vspace macros      chaptergap / sectiongap / tightgap / ...
\end{verbatim}

Chapters reference only these names: no chapter file contains a literal
\code{\textbackslash vspace}, \code{\textbackslash setlength},
\code{\textbackslash renewcommand} or \code{\textbackslash captionsetup} any more.

\section{Known limitations}

\begin{itemize}
  \item End-diaphragm steel tonnage is not produced by the current BOQ calculation
        (it has no diaphragm geometry inputs), so the chart plots it only when the
        key is present. The chart code already supports the category.
  \item The demo PDFs use a synthetic dataset and an inert OpenSees stub; they
        exercise the report generator, not the finite-element solver.
  \item The report generator assembles LaTeX with string builders rather than
        PyLaTeX objects; \code{theme.longtable()} provides the equivalent structured
        API (column spec, head/foot blocks, caption handling) without a dependency
        change.
\end{itemize}

\end{document}
